\documentclass[12pt]{article} 
\renewcommand{\thempfootnote}{\arabic{mpfootnote}}

\usepackage[hidelinks]{hyperref}

\usepackage[utf8]{inputenc}

\usepackage[T1]{fontenc}
\usepackage[ttdefault=true]{AnonymousPro}

\usepackage{multicol}
\setlength{\columnsep}{1cm}

\usepackage{abc}
\usepackage{graphicx}
\usepackage{musicography}

\usepackage[margin=0.5in]{geometry} 
\geometry{letterpaper}

\usepackage{fancyhdr} 
\pagestyle{fancy} % options: empty , plain , fancy
\renewcommand{\headrulewidth}{0pt} % customise the layout...
\lhead{}\chead{}\rhead{}
\lfoot{}\cfoot{\thepage}\rfoot{}

\begin{document}

\thispagestyle{empty}

\begin{abc}[name=westphalia,options={-O= +c -F single.fmt}]
X:1
T:Westphalia Waltz
M:3/4
L:1/8
K:G
DEF|:G4A2|B4e2|(d6|d3) DEF|G4 A2|B4c2| (A4 A)B|A4G2|F4G2|
A4g2|(f4 f)g|f4 e2|d3 ^cde|1 d2^c2=c2|(B4 B)c|B3 DEF:|2 d2F2A2|G6|G2B2d2|
|:g6|f6|e2f2g2|B6|(3BcB A2G2|(3BcB A2G2|c3 Bcd|c6|e6|e6| 
e2f2g2|f4e2|d3 ^cde|1 d2^c2=c2|(B4 B)c|B4d2:|2 d2F2A2|G6|G3|
\end{abc}
\vskip 0.25in
\noindent\makebox[\textwidth][c]{%
  \begin{minipage}{7.0in}
    \setlength{\parindent}{1.5em}

    The western swing waltz, \textbf{Westphalia Waltz}, is associated
    with Texas fiddler, Cotton Collins. Legend has it, in 1946 while
    performing his so-called, `No Name Waltz' in Westphalia Hall, in
    Westphalia, TX, it was suggested Collins name the tune `Westphalia
    Waltz' in honor of the
    town.\footnote{\url{https://secondhandsongs.com/work/235180/all} \\
      \url{https://tunearch.org/wiki/Annotation:Westphalia_Waltz} \\
      \url{http://www.joeweed.com/documentaries/westphalia_waltz.htm}}
    The tune itself predates Collins in that the melody originated
    from the Polish folk tunes, Wszystkie Rybki and Pytala Sie Pani
    (among other names).\footnote{2 According to Joe Weed, ```Pytala
      Sie Pani' was a unifying and bawdy favorite that the overworked,
      underpaid, ostracized and homesick Polish-Americans sang to
      forget the Great Depression.'' According to Google Translate,
      Stanislaw Gresiuk's version of Wszystkie Rybki is about fish and
      drinking alcohol. (See lyrics in video notes of Joe Weed's
      site.)} Various tune origin stories with no solid documentation
    include Collins hearing the song while serving in the US Army in
    Germany during WWII, Collins hearing the song performed in
    Bremond, TX a nearby town to Westphalia with a large Polish
    community, and Collins hearing a recording of the tune as it was
    recorded in America several times during the 1930s.\footnote{secondhandsongs, tunearch} Cotton
    Collins recorded the tune with his band, Lone Star Playboys in
    1946, but it wasn't until Hank Thompson's 1955 recording that the
    tune caught on, eventually becoming extremely popular in the
    Missouri Ozark region as a fiddle contest tune.\footnote{tunearch, Joe Weed}

    ---Kitty Steetle    
  \end{minipage}
}


\end{document}
